\chapter{研究对象、读者与证据方法}
\label{chap:scope-evidence-method}

\section{本书所称“具身智能”}

本书把具身智能定义为：\textbf{通过可观测、可执行的身体与环境形成持续闭环，并在任务、物理和安全约束下感知、估计、决策、行动和更新的智能系统}。这个定义是为比较技术路线而设的操作定义，不替代 ISO 8373 等标准术语\cite{iso8373_2021}。它包含五个必要条件：有明确身体和作用域；有传感与执行接口；动作会改变后续观测；系统针对任务选择行为；失败由物理、任务或安全判据定义。

因此，离线识图模型不是具身系统；把 VLM 接到机械臂上也不自动构成可靠具身智能。反过来，不使用大模型的移动机器人仍可具备闭环自主性。本书关心的不是“是否含 AI”这一标签，而是系统如何在身体和环境约束下形成可复核能力。

\begin{table}[H]
\begingroup
\hbadness=10000
\centering
\caption{研究对象与相邻领域的边界}
\label{tab:scope-boundaries}
\begin{tabularx}{\textwidth}{p{0.16\textwidth}YYY}
\toprule
对象 & 与本书交集 & 本书主讲 & 不在主线的内容 \\
\midrule
机器人学 & 本体、感知、规划、控制 & 学习与基础模型如何进入闭环 & 纯机构设计手册 \\
工业自动化 & 设备、工艺、实时与安全 & 自主策略怎样接入生产系统 & 全套 PLC 编程教程 \\
软件智能体 & 规划、记忆、工具调用 & 动作经身体作用于物理世界 & 纯网页/办公自动化 \\
自动驾驶 & 多传感器、规划、安全 & 可迁移的系统方法 & 道路法规与车辆工程全景 \\
多模态大模型 & 视觉、语言、序列建模 & 与动作、状态和反馈的接口 & 通用聊天模型训练大全 \\
\bottomrule
\end{tabularx}
\endgroup
\end{table}

默认读者具备线性代数、概率和 Python 基础，但未系统学习机器人学。第 4–13 章建立数学、控制和学习基础；第 14–25 章进入多模态、系统架构与 VLA；第 26–30 章讨论数据、评测、部署和安全。产业部分只回答技术交付、供应链和规模化如何受这些约束影响。

\section{本书回答什么，不回答什么}

技术主线回答四类问题：概念和公式如何定义；算法如何训练或求解；代码和系统接口如何实现；证据在哪些任务、硬件和条件下成立。产业主线只在可核事实基础上讨论交付、成本、运维、供应链和政策约束。对于尚无公开证据的商业能力，本书记录“不知道”或“证据不足”，不以融资、演示视频或宣传语补齐结论。

\begingroup\hbadness=10000
本书不是穷尽全部数据库的系统综述，也不进行可合并效应量的元分析（meta-analysis）。PRISMA 2020 强调透明报告检索、筛选和结果\cite{page2021prisma}；本书借鉴其可追溯原则，采用问题导向的教材型技术综述：核心论断优先追溯原始论文、标准、官方代码和可复现实验，并保留检索日期、版本和排除理由。由于机器人任务、硬件和成功判据高度异质，许多结果只能做条件化比较，不能强行合并成统一排名。
\endgroup

\section{证据等级不是一条万能排行榜}

证据强度取决于论断类型。算法机制应看原始论文和源码；接口行为应看对应版本官方文档；标准状态应看标准组织；商业交付应看客户、部署和审计材料。二手综述适合建立广度，不能替代核心模型的原始实验。

\begin{table}[H]
\centering
\caption{论断类型与首选证据}
\label{tab:claim-evidence-map}
\begin{tabularx}{\textwidth}{p{0.20\textwidth}YYp{0.22\textwidth}}
\toprule
论断 & 首选来源 & 必核字段 & 常见误用 \\
\midrule
原理/训练目标 & 原始论文、补充材料 & 假设、公式、数据、消融 & 用新闻稿解释机制 \\
实现/接口 & 官方仓库、固定 commit、文档 & 文件、符号、版本、许可证 & 链接默认分支 \\
性能/SOTA & 原始结果、独立复现 & 任务、硬件、分母、区间 & 跨 benchmark 排名 \\
标准/政策 & 发布机构正式文本 & 编号、版本、状态、范围 & 引用已撤销版本 \\
部署/商业 & 客户验证、运行记录、审计 & 地点、时间、暴露量、介入 & 把试点写成规模交付 \\
图像/代码转载 & 权利页、许可证、作者许可 & 资产、版本、署名、改动 & 把公开可见当可复制 \\
\bottomrule
\end{tabularx}
\end{table}

一个核心论断 $c$ 的最小证据记录写为
\begin{equation}
E(c)=(s,v,l,a,t,r),
\label{eq:claim-evidence-record}
\end{equation}
其中 $s$ 是来源，$v$ 是版本或日期，$l$ 是页码/图号/代码行等定位，$a$ 是适用条件，$t$ 是论断类型，$r$ 是权利和复用状态。只有 URL 而没有定位和版本，难以支撑会变化的实现事实。

\begin{figure}[H]
\centering
\begin{tabular}{c@{\ $\rightarrow$\ }c@{\ $\rightarrow$\ }c@{\ $\rightarrow$\ }c}
\fbox{论断与边界} & \fbox{原始来源/版本} & \fbox{公式、代码或实验} & \fbox{复核与评价}
\end{tabular}
\caption{本书的最小可追溯链。评价只能位于机制和证据之后。}
\label{fig:claim-traceability}
\end{figure}

\section{事实、推断与评价分层}

“论文报告在某任务成功 $k/n$ 次”是事实；“差异可能来自动作频率和复位协议”是推断；“该证据不足以支持通用部署”是评价。三者可以出现在同一节，但不能伪装成同一种确定性。事实使用明确引用；推断列出依据和替代解释；评价说明判断标准。涉及未来、未公开实现或商业数字时，标注时间和不确定性。

\begin{lstlisting}[caption={章节证据包的最低检查；这是写作流程伪代码，不是检索算法}]
for claim in chapter.core_claims:
    require claim.type in {fact, inference, evaluation}
    if claim.type == fact:
        require primary_source(claim) and version_and_locator(claim)
    if claim.concerns_number_or_capability:
        require task_hardware_denominator_and_conditions(claim)
for asset in chapter.figures_tables_code:
    require provenance(asset) and reuse_status(asset)
reject bare_url_as_academic_evidence()
\end{lstlisting}

\section{知识小结与方法判断}

本书的对象是受身体、任务、反馈和安全共同约束的闭环系统。证据按论断匹配来源，并记录版本、定位、适用条件和权利状态。方法判断是：综述的可信度不由参考文献数量单独决定，而由关键论断能否沿“来源—机制/实现—实验—边界”回溯决定；无法回溯的确定语气应被删除或降级。
